\chapter{Vazba praxe na studium}
\label{chap:studies_connection}

Odborná praxe mi umožnila konfrontovat teoretické znalosti získané studiem na Fakultě elektrotechniky a informatiky VŠB-TUO s reálnými požadavky softwarového vývoje. V mnoha oblastech jsem mohl přímo navázat na absolvované předměty, zatímco jiné domény vyžadovaly intenzivní samostudium.

\section{Uplatnění teoretických znalostí}

Nejvýznamnějším přínosem studia pro tuto praxi byla znalost objektově orientovaného programování a databázových systémů. Konkrétně jsem využil znalosti z následujících předmětů:

\begin{itemize}
    \item \textbf{Programování v C\# 1 a 2:} Znalost syntaxe jazyka C\#, typového systému a asynchronního programování byla nezbytná pro vývoj backendové části (Service Fabric Actor) a psaní automatizovaných testů.
    \item \textbf{Databázové systémy 1 a 2:} Při návrhu nových tabulek pro modul Importér a optimalizaci dotazů jsem uplatnil principy normalizace dat, tvorby indexů a práci s cizími klíči.
    \item \textbf{Funkcionální programování (FPR):} Ačkoliv je většina projektu psána v imperativních jazycích, principy rekurze probírané v tomto předmětu byly klíčové při implementaci hierarchického stromu kontextů (\textit{Context Tree}). Zobrazení vnořených struktur (např. Subjekt $\rightarrow$ Adresa $\rightarrow$ Kontakt) vyžadovalo rekurzivní průchod daty, což je koncept, který jsem si osvojil právě v FPR.
    \item \textbf{Správa softwarových projektů (SSP):} Základy verzování kódu pomocí systému Git, které byly představeny v tomto předmětu, jsem využíval denně při práci s větvemi (branches) a správě zdrojových kódů.
\end{itemize}